\section{Lie Groups}

\subsection{Basic definitions}

\bdefn

    A {\emphcolor Lie group} is a smooth manifold $G$ without boundary, along with smooth maps
    $$ m\colon G\times G\to G,\qquad i\colon G\to G $$
    turning $G$ into a group with multiplicatio given by $m$ and inversion given by $i$.

\edefn

\bnote

    A Lie group is a special type of {\emphcolor topological group}, which is a group which is also a topological space making
    multiplication and inversion into continuous functions.

\enote

Note that $G$ is a Lie group iff the map $(g,h)\mapsto gh^{-1}$ is smooth.
This map is a composition of multiplication and inversion, so if $G$ is a Lie group it is smooth, and conversely we note that it being smooth
implies $i$ is smooth, and then that $m$ is smooth.

For a Lie group $G$ and $g\in G$, define {\emphcolor left translation} and {\emphcolor right translation} by
$$ L_g,R_g\colon G\to G,\qquad L_g(h) = gh,\quad R_g(h) = hg . $$
These are smooth maps, as they are simply compositions of $m$ with a suitable inclusion $G\to G\times G$.
Moreso, they are diffeomorphisms: the inverse of $L_g$ is $L_{g^{-1}}$ and similar for right translations.

\bexam

    Euclidean space $\bR^n$ is a Lie group under addition, clearly.

\eexam

\bexam

    The general linear group $\gl_n(\bR)$ is an open submanifold of the Euclidean space $\mat_n(\bR)$.
    Furthermore it is a Lie group, with multiplication and inversion given by usual multiplication and inversion of matrices.
    These are smooth, as the coordinates of multiplication are polynomial in their inputs and inversion can be written as
    $$ A^{-1} = 1/\det(A)\cdot{\rm adj}(A) , $$
    the product of smooth functions (the coordinates of the adjugate are determinants of minors).

    Denote by $\gl_n^+(\bR)$ the subgroup of $\gl_n(\bR)$ with positive determinant.
    This is an open subspace of $\gl_n(\bR)$ (the preimage of $(0,\infty)$) and the restriction of multiplication and inversion to it
    is still smooth (this is true for embedded submanifolds in general).

\eexam

Continuing the above example, if $G$ is a Lie group and $H\subseteq G$ is a subgroup which is also an open set, then it is also a Lie group.

\bexam

    The complex general linear group $\gl_n(\bC)$ is the group of invertible matrices in $\mat_n(\bC)$.
    This is an open subset of $\mat_n(\bC)\cong\bR^{2n^2}$ and thus a $2n^2$-dimensional Lie group.

\eexam

\bdefn

    If $G$ and $H$ are Lie groups, a {\emphcolor Lie group morphism} $G\to H$ is a smooth map which is also a group morphism.
    A {\emphcolor Lie group isomorphism} is a Lie group morphism which is also a diffeomorphism (and in particular a group isomorphism).

\edefn

\bexam

    Consider $\bR$ as a Lie group under addition, and $\bR^\times$ a Lie group under multiplication.
    (We need not show that these are Lie groups again, $\bR$ is Euclidean and $\bR^\times=\gl_1(\bR)$.)
    Then the exponential map, $\exp\colon\bR\to\bR^\times$, is a Lie group morphism.
    Its image is $\bR^+$, and it has an inverse $\log\colon\bR^+\to\bR$, so that $\bR^+$ (with multiplication) and $\bR$ (with addition)
    are isomorphic Lie groups.

\eexam

\bexam

    The determinant $\det\colon\gl_n(\bC)\to\bR^\times$ is a Lie group morphism.

\eexam

Most importantly, if $G$ is a Lie group and $g\in G$, then {\emphcolor conjugation} by $g$,
$$ C_g\colon G\to G,\qquad C_g(h) = ghg^{-1} $$
is a Lie group morphism.
It is clearly smooth, as a composition of multiplication and inversion, and it is an isomorphism with inverse $C_{g^{-1}}$.
Recall that $H\subseteq G$ is {\emphcolor normal} if $C_g(H)=H$ for all $g\in G$.

\bthrm

    Every Lie group morphism has constant rank.

\ethrm

\bproof

    Let $F\colon G\to H$ be a Lie group morphism, with $e,\tilde e$ the respective units of the groups.
    Now let $g_0\in G$ and notice that
    $$ F(L_{g_0}(g)) = F(g_0g) = F(g_0)F(g) = L_{F(g_0)}F(g) $$
    so that $F\circ L_{g_0}=L_{F(g_0)}\circ F$.
    Thus we get that
    $$ dF_{g_0}\cdot d(L_{g_0})_e = d(L_{F(g_0)})_{\tilde e}\cdot dF_e . $$
    Since left translations are diffeomorphism, they do not change rank, and thus $dF_{g_0}$ and $dF_e$ have the same rank, meaning
    $F$ has constant rank.
    \qed

\eproof

\bcoro

    A Lie group morphism is an isomorphism iff it is bijective.

\ecoro

\bproof

    A bijective Lie group morphism is a bijective immersion (by constant rank), thus a diffeomorphism.
    \qed

\eproof

\subsection{Lie subgroups}

\bdefn

    Let $G$ be a Lie group.
    A {\emphcolor Lie subgroup} of $G$ is a subgroup $H\leq G$ endowed with a smooth structure making it a Lie group and an immersed submanifold.

\edefn

\bprop

    Let $G$ be a Lie group, and $H\leq G$ a subgroup which is an embedded submanifold.
    Then $H$ is also a Lie group, and thus a Lie subgroup of $G$.

\eprop

\bproof

    Since $H$ is embedded, multiplication $G\times G\to G$ restricts both on its domain (for all submanifolds) and its codomain (as $H$ is embedded),
    to give $H$'s multiplication $H\times H\to H$.
    Same for inversion.
    \qed

\eproof

\blemm

    Let $G$ be a Lie group and $H\subseteq G$ an open subgroup.
    Then $H$ is an embedded Lie subgroup.
    Moreso, $H$ is also closed, and thus a union of connected components of $G$.

\elemm

\bproof

    By the above proposition, $H$ is an embedded Lie subgroup.
    Furthermore, each coset $gH$ is open in $G$ as the image of $H$ under the diffeomorphism $L_g$.
    Then $G-H=\bigcup_{g\in G-H}gH$ is open, and so $H$ is closed.
    Clopen sets are unions of connected components.
    \qed

\eproof

\bprop

    Let $G$ be a Lie group, and $W\subseteq G$ an open, then
    \benum
        \item $W$ generates an open subgroup of $G$,
        \item if $W$ is a connected neighborhood of the identity, it generates a connected open subgroup of $G$,
        \item if $G$ is connected, then $W$ generates $G$.
    \eenum

\eprop

\bproof

    Denote for arbitrary sets $A,B$,
    $$ AB = \set{ab}[a\in A,b\in B],\qquad A^{-1} = \set{a^{-1}}[a\in A] . $$
    For each $k>0$ define by $W_k$ the set of all elements in $G$ that can be written as a product of $k$ or fewer elements of $W\cup W^{-1}$.
    Then $W$'s generated subgroup is $\bigcup_{k>0}W_k$.

    Note that $W_1=W\cup W^{-1}$ (we do not allow empty products, defined to be the identity, so if $e\notin W$ then it is not in $W_1$, but is
    in $W_2$).
    This is open, as $W^{-1}=i(W)$ is the diffeomorphic image of an open set.
    And inductively we have
    $$ W_k = W_1W_{k-1} = \bigcup_{g\in W_1}L_g(W_{k-1}) $$
    is open.
    Thus the subgroup generated by $W$ is a union of open sets and thus open.

    If $W$ is a connected neighborhood of $e$, then so is $W^{-1}$ (as it is the diffeomorphic image of a connected set, and contains $e$).
    Then $W_1=W\cup W^{-1}$ is also a connected neighborhood of $e$, and inductively so are $W_k=m(W_1\times W_{k-1})$, thus so is their union.
    This is because $W_1\times W_{k-1}$ is connected, and the continuous image of a connected space is connected, and the union of connected
    spaces with nontrivial intersection is connected (each $W_k$ contains $e$).

    If $G$ is connected, then the subgroup $W$ generates is open by point (1), and thus closed.
    Only trivial subspaces of a connected space are clopen, so $W$ generates $G$.
    \qed

\eproof

For a Lie group $G$, the connected component containing the identity is called the {\emphcolor identity component}.

\bprop

    Let $G_0$ be the identity component of $G$.
    Then $G_0$ is a normal subgroup of $G$, and it is the only connected open subgroup.
    Every connected component of $G$ is diffeomorphic to $G_0$.

\eprop

\bproof

    Any connected open neighborhood of the identity generates $G_0$ by the above proposition.
    It is normal since $C_g(G_0)$ is a connected neighborhood of the identity, and thus contained in $G_0$ for all $g$.
    And $L_g$ maps $G_0$ diffeomorphically to the connected component of $g\in G$.
    \qed

\eproof

\bprop

    Let $F\colon G\to H$ be a Lie group morphism, then $\ker F$ is a properly embedded Lie subgroup of $G$, whose codimension is equal to
    $F$'s rank.

\eprop

\bproof

    $\ker F$ is the level set $F^{-1}e$, and since $F$ has constant rank, it is a properly embedded submanifold of $G$.
    We already showed that subgroups which are embedded submanifolds of Lie groups are Lie subgroups.
    \qed

\eproof

\bprop

    If $F\colon G\to H$ is an injective Lie group morphism, its image has a unique smooth manifold structure turning it into a Lie
    subgroup of $H$ with $F$ an isomorphism onto its image.

\eprop

\bproof

    $F$ is injective of constant rank, and thus an immersion.
    Thus $F(G)$ has a unique smooth manifold structure turning it into an immersed submanifold of $H$ diffeomorphic to $G$.
    Since it is diffeomorphic to $G$, it is also a Lie group and a subgroup (for obvious algebraic reasons).
    Thus it is a Lie subgroup.
    \qed

\eproof

\bexam

    The set $\Sl_n(\bR)$ of $n\times n$ real matrices of determinant $1$ is a properly embedded Lie subgroup of $\gl_n(\bR)$, as the kernel
    of the determinant.
    The rank of the determinant map is $1$, so $\Sl_n(\bR)$ has dimension $n^2-1$.

\eexam

\bexam

    The set $\Sl_n(\bC)$ of $n\times n$ complex matrices of determinant $1$ is the kernel of the determinant map to $\bC^\times$.
    This is a surjective map, and thus has rank $2$ (since $\bC^\times$ does).
    Thus $\Sl_n(\bC)$ is a properly embedded submanifold of $\gl_n(\bC)$ of dimension $2n^2-2$.

\eexam

\bthrm

    Let $H\subseteq G$ be a Lie subgroup.
    Then $H$ is closed iff it is embedded.

\ethrm

\bproof

    If $H$ is embedded, let $g$ be a limit point of $H$ (recall that in Hausdorff spaces, a set is closed iff it contains its limit points).
    Then there is a sequence of elements $h_i\in H$ converging to $g$.
    Let $(U,\phi)$ be a chart for $G$ containing the identity such that $\iota$'s coordinate representation is concatenation with $0$.
    Then we can take $W$ to be a smaller neighborhood so that $\cl W\subseteq U$.
    We can take a neighborhood $V$ of $e$ such that $g_1g_2^{-1}\in W$ whenever $g_1,g_2\in V$.
    Consider the smooth map $G\times G\to G$ by $g_1,g_2\mapsto g_1g_2^{-1}$ and then take $W$'s preimage, which is open and thus contains
    a neighborhood $V_1\times V_2$ of the identity, and we can shrink them to give $V_1=V_2=V$.

    Since $h_ig^{-1}\to e$, we can assume that $h_ig^{-1}\in V$ for all $i$ (since eventually this is true).
    Thus
    $$ h_ih_j^{-1} = (h_ig^{-1})(h_jg^{-1})^{-1} \in W $$
    for all $i,j$.
    Fixing $j$ and letting $i\to\infty$, we see that $h_ih_j^{-1}\to gh_j^{-1}\in\cl W\subseteq U$.
    By our assumption, $H\cap U$ is closed and thus $gh_j^{-1}\in H$, so $g\in H$.

    Conversely let $H$ be a closed Lie subgroup with $m=\dim H,n=\dim G$.
    If $m=n$ then it must be embedded as full-rank immersed submanifolds are embedded.
    So assume $m<n$.

    It is sufficient to show that for some $h_1\in H$ there is a neighborhood $h_1\in U_1\subseteq G$ such that $H\cap U_1$ is an embedded submanifold
    of $U_1$.
    For any other $h\in H$, consider the right translation $L_{hh_1^{-1}}\colon G\to G$ which maps $H$ to $H$ and $U_1$ to a neighborhood
    $U'_1$ of $h$.
    Then $U'_1\cap H$ is embedded in $U'_1$.
    So for every $h\in H$ there is a neighborhood $U\subseteq G$ such that $U\cap H$ is embeddable in $U$.
    Then by the local slice criterion, $H$ is embeddable in $G$.

    Every immersion is locally an embedding, so let $h_1\in V\subseteq G$ such that $\iota|_V\colon V\cap H\to V$ is an embedding.

\eproof

\TODO{Above}

\subsection{Lie group actions}

Recall that a (left) group action of a group $G$ on a set $X$ is a map $\theta\colon G\times X\to X$ written by juxtaposition satisfying
$$ g_1(g_2x) = (g_1g_2)x,\qquad ex = x . $$
A right group action is defined similarly.

\bdefn

    Let $G$ be a topological group and $X$ a topological space.
    A {\emphcolor continuous (left) group action} of $G$ on $X$ is a group action $\theta\colon G\times X\to X$ which is continuous.
    A {\emphcolor smooth group action} is a group action of a Lie group on a smooth manifold which is smooth.

\edefn

There is a duality between left and right group actions: a left group action defines a right group action by $xg=g^{-1}x$ and vice versa.

Note that regular group actions can be equivalently characterized as group morphisms $G\to\sym(X)$ where $\sym(X)$ is the set of permutations of $X$.
For continuous group actions, since the category of topological spaces doesn't in general have a right adjoint to products, there isn't in general
a topological space like $\sym(X)$.

But $\cdot\times X$ does have a right adjoint if $X$ is locally compact and Hausdorff.
Its right adjoint is $Z^X$, the set of continuous maps $X\to Z$, given the compact-open topology: for compact $K\subseteq X$ and open $U\subseteq Z$,
we add base sets $V(K,U)$, which are continuous maps $f\colon X\to Z$ where $f(K)\subseteq U$.
Let $\aut(X)$ be the topological group of automorphisms of $X$, viewed as a subspace of $X^X$.
In this case, a continuous group action is the same as a continuous group morphism
$$ G \to \aut(X) . $$
But when $X=M$ is a smooth manifold, we then need to give $\aut(M)$ a Lie group structure.

For a group action $\theta\colon G\times M\to M$, for $g\in G$ we write the map $\theta_g\colon M\to M$.
These are smooth maps, as $\theta$ is, and satisfy
$$ \theta_g\circ\theta_h = \theta_{gh},\qquad \theta_e=\id_M . $$
Thus $\theta_g$ is a diffeomorphism with inverse $\theta_{g^{-1}}$.

\bdefn

    For a group action of $G$ on $M$, define
    \benum
        \item the {\emphcolor orbit} of $p\in M$ to be $G\cdot p=\set{gp}[g\in G]$,
        \item the {\emphcolor stabilizer} of $p\in M$ to be $G_p=\set{g\in G}[gp=p]$ (also called the {\emphcolor isotropy group}).
    \eenum
    The isotropy group is indeed a subgroup of $G$.
    The action is
    \benum
        \item {\emphcolor transitive} if for every $p,q\in M$ there is a $g\in G$ such that $gp=q$, equivalently $M$ has a single orbit.
        \item {\emphcolor free} if every isotropy group is trivial: if $gp=p$ then $g=e$.
    \eenum

\edefn

\bexam

    If $G$ is a Lie group and $M$ a smooth manifold, the {\emphcolor trivial action} of $G$ on $M$ is $gp=p$ for all $g\in G,p\in M$.

\eexam

\bexam

    $\gl_n(\bR)$ has a natural action on $\bR^n$ by matrix multiplication: $(A,v)\mapsto Av$.
    This is clearly a group action, and smooth since its components are all polynomially dependent on the inputs.
    $\gl_n(\bR)$ almost acts transitively: for every $v,u\in\bR^n$ which are nonzero there is an invertible matrix $A$ such that $Av=u$.
    So this action has two orbits: $\bR^n-0$ and $0$.

\eexam

If $H\subseteq G$ is a Lie subgroup, and $G$ acts on $M$, then we can restrict the action to $H$, and it remains smooth.
This is because $H\times M$ is an immersed submanifold of $G\times M$.

\bexam

    Left translation defines a smooth action of a Lie group on itself (right translation defines a right smooth action).
    This is transitive, as $g_2g_1^{-1}$ maps $g_1$ to $g_2$, and it is free.
    The restriction of this to a Lie subgroup remains free and smooth, but not necessarily transitive.

\eexam

\bexam

    Conjugation forms a smooth group action of a Lie group on itself.

\eexam

\bdefn

    If $G$ acts smoothly on two manifold $M$ and $N$, a {\emphcolor $G$-equivariant} map $F\colon M\to N$ is a smooth map which is equivariant
    in the usual way:
    $$ F(g\cdot p) = g\cdot F(p) $$
    for all $g\in G,p\in M$.

\edefn

If the actions of $G$ on $M$ and $N$ are given by $\theta,\phi$ respectively, this is the same as saying that the following diagram commutes for
all $g\in G$:

\centerline{\drawdiagram{
    $M$&$N$\cr
    $M$&$N$\cr
}{
    \diagarrow{from={1,1}, to={1,2}, text=$F$, y distance=-.25cm}
    \diagarrow{from={2,1}, to={2,2}, text=$F$, y distance=.25cm}
    \diagarrow{from={1,1}, to={2,1}, text=$\theta_g$, x distance=-.25cm}
    \diagarrow{from={1,2}, to={2,2}, text=$\phi_g$, y distance=.25cm}
}}

\bthrm

    Let $G$ act smoothly on $M$ and $N$ such that it acts transitively on $M$.
    If $F\colon M\to N$ is $G$-equivariant, it has constant rank.

\ethrm

\bproof

    For $p,q\in M$ let $g\in G$ be such that $gp=q$ (by transitivity).
    Then we have $\phi_g\circ F=F\circ\theta_g$ and thus
    $$ d(\phi_g)_{F(p)}\cdot dF_p = dF_q\cdot d(\theta_g)_p , $$
    and since the actions are diffeomorphisms, we have that $dF_p$ and $dF_q$ have the same rank, as desired.
    \qed

\eproof

For a smooth action $\theta\colon G\times M\to M$, consider the smooth maps for $p\in M$ defined by
$$ \theta^{(p)}\colon G\to M,\qquad \theta^{(p)}(g) = gp . $$
This is called the {\emphcolor orbit map}, as its image is the orbit $G\cdot p$.
The preimage $(\theta^{(p)})^{-1}(p)$ is the isotropy group $G_p$.

\bcoro

    For any smooth action $\theta$ and $p\in M$, $\theta^{(p)}\colon G\to M$ is smooth and has constant rank.
    Thus the isotropy group $G_p=(\theta^{(p)})^{-1}(p)$ is a properly embedded Lie subgroup of $G$.
    If $G_p=\set e$, then $\theta^{(p)}$ is injective and thus a smooth immersion, so $G\cdot p$ is an immersed submanifold of $M$.

\ecoro

\bproof

    $\theta^{(p)}$ is equivariant relative to the group action of $G$ on itself and $G$ on $M$:
    $$ \theta^{(p)}(g_1g_2) = g_1g_2p = g_1\theta^{(p)}(g_2) , $$
    and since $G$ acts transitively on above, this has constant rank as above.
    Thus $G_p$ is a level set of a constant rank map, and thus a properly embedded smooth manifold, and thus a Lie subgroup.
    If $G_p$ is trivial, then $\theta^{(p)}$ is clearly injective, and thus a smooth immersion, and the result follows.
    \qed

\eproof

\bexam

    A matrix $A\in\mat_n(\bR)$ is {\emphcolor orthogonal} if it preserves the Euclidean dot product:
    $$ (Ax)\cdot(Ay) = x\cdot y , $$
    or equivalently its columns are orthonormal, i.e.\ $A^\top A=I_n$.
    Let $\orth(n)$ be the set of all orthogonal real $n\times n$ matrices.
    So define $\Phi\colon\gl_n(\bR)\to\mat_n(\bR)$ by $\Phi(A)=A^\top A$, so that $\orth(n)$ is the level set $\Phi^{-1}I_n$.

    To show that $\orth(n)$ is an embedded submanifold of the general linear group, we will show that it is equivariant for some group action.
    There is the clear transitive smooth action of $\gl_n(\bR)$ on itself, and define its action on $\mat_n(\bR)$ by a sort of conjugation:
    $$ A\cdot X = AXA^\top . $$
    Then $\Phi$ is equivariant:
    $$ \Phi(AB) = (AB)(AB)^\top = ABB^\top A^\top = A\cdot\phi(B) , $$
    as desired.

    $\orth(n)$ is compact, as it is closed and bounded in $\mat_n(\bR)$: closed as the preimage of a closed set, and bounded since norms of
    orthogonal matrices are
    $$ \norm A^2 = \trace(A^\top A) = \trace(I_n) = n , $$
    so bound by $\sqrt n$.

    To compute $\orth(n)$'s dimension we need to compute $\Phi$'s rank, which we can do at any matrix of our choosing, so take $I_n$.
    To do this, we compute $d\Phi$'s image on tangent vectors, viewed as velocity curves.
    So for $B\in\mat_n(\bR)$ we consider the curve $\gamma\colon(-\epsilon,\epsilon)\to\gl_n(\bR)$ by $\gamma(t)=I_n+tB$.
    Since $T_{I_n}\gl_n(\bR)\cong\mat_n(\bR)$, all curves are uniquely of this form.
    Then
    $$ d\Phi_{I_n}(B) = \eval{\frac d{dt}}_{t=0}\Phi\circ\gamma(t) = \eval{\frac d{dt}}_{t=0}(I_n+tB)^\top(I_n+tB) = B^\top + B . $$
    Thus $d\Phi$'s image is precisely the symmetric matrices, which has dimension $n+n(n-1)/2=n(n+1)/2$.
    And so $\orth(n)$'s dimension is $n^2-n(n+1)/2=n(n-1)/2$.

\eexam

\bexam

    The {\emphcolor special orthogonal group} of degree $n$ is the set of real orthogonal matrices of determinant $1$.
    That is, $\sorth(n)=\orth(n)\cap\sl_n(\bR)$.
    Notice that for $A\in\orth(n)$ we have
    $$ 1 = \det I_n = \det(A^\top A) = \det(A)^2 , $$
    so that $\sorth(n)$ is the set of orthogonal matrices of positive determinant, thus an open submanifold of $\orth(n)$.

\eexam

